% BEGIN VOL07-EXPANSION VII13
\section{Supplementary worked examples}
\label{sec:vii13-pedagogy}

\begin{example}[Frenet frame of a circle]\label{ex:vii13-ped-01}
For the unit-speed circle \(\gamma(s)=R(\cos(s/R),\sin(s/R),0)\),
\[
\kappa=\frac1R,\qquad \tau=0.
\]
The principal normal points toward the center, and the binormal is constant.  Curvature measures the turning of the tangent, while torsion vanishes because the curve is planar.
\end{example}

\begin{example}[Curvature and torsion of a helix]\label{ex:vii13-ped-02}
For \(\gamma(t)=(a\cos t,a\sin t,bt)\), one finds
\[
\kappa=\frac{a}{a^2+b^2},\qquad
\tau=\frac{b}{a^2+b^2}.
\]
Both are constant.  The ratio \(\tau/\kappa=b/a\) records the pitch relative to the radius.
\end{example}

\begin{example}[Why the Frenet frame can fail]\label{ex:vii13-ped-03}
If a regular curve has \(\kappa=0\) at a point, then \(T'=0\) there and the principal normal \(N=T'/\|T'\|\) is undefined.  The Frenet frame therefore requires more than regularity; it requires nonzero curvature on the interval under discussion.
\end{example}

\section{Graded supplementary exercises}
The following exercises add a deliberate progression from concrete calculation to structural proof, hypothesis testing, application, and synthesis.

\subsection*{Standard computations and constructions}

\begin{exercise}[Circle curvature]\label{exr:vii13-ped-01}
Compute the curvature of a circle of radius \(R\).
\end{exercise}
\begin{hint}
Use a unit-speed parametrization.
\end{hint}
\begin{solution}
The unit tangent changes at rate \(1/R\), so \(\kappa=1/R\).
\end{solution}

\begin{exercise}[Planar torsion]\label{exr:vii13-ped-02}
Compute the torsion of a planar circle viewed in \(\mathbb R^3\).
\end{exercise}
\begin{hint}
The binormal is constant.
\end{hint}
\begin{solution}
\(\tau=0\).
\end{solution}

\begin{exercise}[Helix invariants]\label{exr:vii13-ped-03}
For \(\gamma(t)=(3\cos t,3\sin t,4t)\), compute \(\kappa\) and \(\tau\).
\end{exercise}
\begin{hint}
Use the standard helix formulas.
\end{hint}
\begin{solution}
\(\kappa=3/25\) and \(\tau=4/25\).
\end{solution}

\begin{exercise}[Osculating circle]\label{exr:vii13-ped-04}
Find the radius of the osculating circle of a curve with curvature \(\kappa=5\).
\end{exercise}
\begin{hint}
Radius of curvature is \(1/\kappa\).
\end{hint}
\begin{solution}
The radius is \(1/5\).
\end{solution}

\begin{exercise}[Straight line]\label{exr:vii13-ped-05}
Compute curvature for a unit-speed straight line.
\end{exercise}
\begin{hint}
Its tangent is constant.
\end{hint}
\begin{solution}
\(\kappa=0\).
\end{solution}

\subsection*{Proofs}

\begin{exercise}[Frenet orthogonality]\label{exr:vii13-ped-06}
Prove \(T,N,B\) form an orthonormal frame when \(\kappa>0\).
\end{exercise}
\begin{hint}
Use \(T\cdot T=1\) and \(B=T\times N\).
\end{hint}
\begin{solution}
\(T'\perp T\), hence \(N\perp T\) and has unit length; \(B\) is the unit cross product orthogonal to both.
\end{solution}

\begin{exercise}[Planar curves have zero torsion]\label{exr:vii13-ped-07}
Prove a curve contained in an affine plane has \(\tau=0\) wherever its Frenet frame exists.
\end{exercise}
\begin{hint}
The binormal is a fixed unit normal to the plane.
\end{hint}
\begin{solution}
Since \(B\) is constant, \(B'=0=-\tau N\), so \(\tau=0\).
\end{solution}

\begin{exercise}[Rigid-motion invariance]\label{exr:vii13-ped-08}
Prove curvature is invariant under Euclidean rigid motions.
\end{exercise}
\begin{hint}
Orthogonal maps preserve norms and derivatives.
\end{hint}
\begin{solution}
A rigid motion sends the unit tangent derivative to its orthogonal image, preserving \(\|T'\|\), hence curvature.
\end{solution}

\begin{exercise}[Frenet equations preserve orthonormality]\label{exr:vii13-ped-09}
Explain why the skew-symmetric coefficient matrix in the Frenet equations is forced by orthonormality.
\end{exercise}
\begin{hint}
Differentiate all pairwise inner products.
\end{hint}
\begin{solution}
For an orthonormal moving frame \(F\), differentiating \(F^TF=I\) gives \(F'^TF+F^TF'=0\), so the connection matrix is skew-symmetric.
\end{solution}

\subsection*{Counterexamples and hypothesis tests}

\begin{exercise}[Regularity guarantees a Frenet frame]\label{exr:vii13-ped-10}
Test with a straight line.
\end{exercise}
\begin{hint}
Its curvature is zero.
\end{hint}
\begin{solution}
False.  A straight line is regular but \(N\) and \(B\) are not defined by the Frenet construction.
\end{solution}

\begin{exercise}[Zero torsion means straight line]\label{exr:vii13-ped-11}
Test with a circle.
\end{exercise}
\begin{hint}
A circle is planar.
\end{hint}
\begin{solution}
False.  The circle has zero torsion but nonzero curvature.
\end{solution}

\begin{exercise}[Curvature changes under reparametrization]\label{exr:vii13-ped-12}
Test for regular orientation-preserving reparametrizations.
\end{exercise}
\begin{hint}
Curvature is geometric.
\end{hint}
\begin{solution}
False.  When computed correctly, curvature is invariant under regular reparametrization.
\end{solution}

\subsection*{Applications and investigations}

\begin{exercise}[Turning radius]\label{exr:vii13-ped-13}
Explain the geometric meaning of \(1/\kappa\).
\end{exercise}
\begin{hint}
Compare the curve locally with its osculating circle.
\end{hint}
\begin{solution}
The reciprocal curvature is the radius of the circle with second-order contact in the normal plane when \(\kappa\neq0\).
\end{solution}

\begin{exercise}[Helix pitch]\label{exr:vii13-ped-14}
How do \(\kappa\) and \(\tau\) distinguish helices of different pitch?
\end{exercise}
\begin{hint}
Use \(\tau/\kappa=b/a\).
\end{hint}
\begin{solution}
The ratio measures axial rise relative to radius; larger \(|\tau/\kappa|\) corresponds to a more stretched helix.
\end{solution}

\subsection*{Challenge problems}

\begin{exercise}[Recover a constant-invariant helix]\label{exr:vii13-ped-15}
Assuming constant \(\kappa>0\) and constant \(\tau\), explain why the curve is congruent to a circular helix.
\end{exercise}
\begin{hint}
Use the Frenet ODE with constant coefficients.
\end{hint}
\begin{solution}
The Frenet system has a constant skew-symmetric coefficient matrix; solving it yields a frame rotating about a fixed axis, and integrating \(T\) produces a circular helix, with the circle as the \(\tau=0\) case.
\end{solution}

\begin{exercise}[Torsion sign and orientation]\label{exr:vii13-ped-16}
Explain what happens to torsion when the ambient orientation of \(\mathbb R^3\) is reversed while the curve orientation is fixed.
\end{exercise}
\begin{hint}
The cross product and binormal change sign.
\end{hint}
\begin{solution}
The binormal changes sign, so the signed torsion changes sign; curvature remains unchanged.
\end{solution}

% END VOL07-EXPANSION VII13
